Calculation of the moment of inertia of Earth \hfill[2]

The moment of inertia is
\[ I = \frac{2}{5} m_{\mathrm{E}} R_{\mathrm{E}}^2 \]

Equation of motion for this case
\[ K = I\cdot\varepsilon \]
\hfill[2]

where $\varepsilon$ is angular acceleration and $K$ is moment of force.

Calculation of angular acceleration \hfill[1]

Acceleration is equal to $\varepsilon = K/I$. From the assumption, the
acceleration $\varepsilon$ is constant and we find its numerical value equal
to $6\cdot10^{-22}\,\mathrm{s}^{-2}$. % sic: the minus sign of the exponent is illegible in the scan

Calculation of the change of $\varepsilon$ during assumed period of time % sic: change of omega, not epsilon
\hfill[1]

Multiplying this value by $1.9\cdot10^{16}$\,s [600\,My] we obtain
$\Delta\omega = 1.14\cdot10^{-5}\,\mathrm{s}^{-1}$.

Calculation of the angular frequency in the past \hfill[2]

Using the definition of the angular frequency one can calculate its recent
value $2\pi/86164\,\mathrm{s}^{-1} = 7.29\cdot10^{-5}\,\mathrm{s}^{-1}$.
Subtracting from it calculated above $\Delta\omega$ one gets
$\omega_{\mathrm{past}} = 6.15\cdot10^{-5}\,\mathrm{s}^{-1}$.

Calculation of the day length in the past and number of days in the ancient
year \hfill[2]

Having angular frequency one can get ancient period of Earth rotation.
Dividing sidereal year length by the obtained period one gets number of days
in the ancient year equal to about 420 days.
